\section{Limitations and open questions}\label{sec:outlook}

\paragraph{The results are asymptotic.}
The proofs of Theorems~\ref{thm:planar}--\ref{thm:resonant} do not compute $\mu_0$ or the other constants. They
come from compactness: the widths of the complex domain around the whole sublevel, the bound on the complexified
perturbation there, the modulus of continuity of $B$, and, for Theorem~\ref{thm:resonant}, the thresholds of
Proposition~\ref{thm:partialnek}, which are not explicit either. The theorems say nothing about a given value of
$\mu$, and no statement about a real planetary system follows from them.

\paragraph{Real masses.}
An explicit version would need rigorous lower bounds for the complex widths near the boundary of the sublevel,
where perihelia and mutual distances approach their minima, rigorous bounds for the complexified perturbation
on the whole enlarged slow set, and, for Theorems~\ref{thm:planar} and~\ref{thm:spatial}, control of the small
divisors up to a cutoff. With generic majorants one should expect admissible mass ratios many orders of magnitude
below planetary values. Tailored estimates do better, but remain far from planetary masses: for a planar
Sun--Jupiter--Saturn model near the $5{:}2$ resonance, \cite[\S5.1]{BN20} reaches times comparable to the age of
the Solar System for perturbations of size about $10^{-11}$ in their normalization, while the Jupiter--Sun mass
ratio is about $10^{-3}$. Nor can observations establish a Diophantine condition to all orders, since every
interval of mean motions contains resonances of high order; Theorem~\ref{thm:resonant} removes this hypothesis,
but an explicit version would still have to control resonances up to a finite order.

\paragraph{The Solar System.}
We make no claim about the Solar System. Laskar and Petit find it AMD-unstable, because of the AMD carried by the
giant planets \cite[\S5.4.1]{LP17}, so the hypothesis of our theorems fails for the full eight-planet Newtonian
model, whatever the constants. The model \eqref{eq:H} also omits all non-Newtonian effects.

\paragraph{Open questions.}
\begin{enumerate}
\item \emph{Explicit constants.} Can the proof be made quantitative, perhaps computer-assisted with a finite
  Fourier cutoff and rigorous divisor bounds, for a published model of a system with a large AMD margin?
\item \emph{Exponents.} The exponents of Theorem~\ref{thm:resonant} are those of Nekhoroshev's general theorem, as
  in \cite[Theorem~12.3]{Nek77}. Can they be improved on the sharp sublevel, using the convexity of the Kepler
  Hamiltonian or the structure of specific resonances?
\item \emph{Other thresholds.} What are the analogues of our theorems for the weaker notion of AMD-stability of
  \cite[\S5.3]{LP17}, which allows a collision of the innermost planet with the star?
\item \emph{Longer times.} Does a fixed margin below the strong threshold give stability beyond exponentially long
  times?
\item \emph{Earlier work.} Niederman \cite{Nie96} proved exponential stability results in the planetary problem.
  We have not inspected that paper; a precise comparison of hypotheses, domains and exponents remains to be made.
\end{enumerate}
